% !TeX root = ../papers/nc-p3.tex
\section{順次中心になるパラメータに対する一様な有限判定}
\label{sec:successively-central-cutoff}
\begin{theorem}\label{thm:successively-central-cutoff}
標準有限Frobenius候補の折り畳み環 $B$ を取り、正次数 $d_i$ の斉次な頂点対角元
$g_1,\ldots,g_4$ が $B_{i-1}=B/(g_1,\ldots,g_{i-1})$ で順に中心になるとする。
$B_4$ が有限次元なら次が左右Koszul AS正則性の必要十分条件である。
$L=\sum(d_i-1)$、$b=\max(1,L,d_1,\ldots,d_4)$ とする。
\begin{enumerate}
\item 各始点で $B_4$ の次数 $0,\ldots,L+1$ のHilbert関数が
$\prod_i(1+z+\cdots+z^{d_i-1})$ の係数に等しい。
\item $B$ の各始点の経路次元は $3b$ 次まで標準で、左右の一方の標準Koszul複体は
内部次数 $L+4$ まで完全である。
\end{enumerate}
第一条件があれば有限性も同時に証明される。成立すれば四元は順に両側非零因子である。
四元とも四次なら経路次元は36次、複体は16次までで足りる。
各元は前の元を割った環でだけ中心でよく、元の環で四元が生成する部分代数の表示は要らない。
\end{theorem}
\begin{proof}
まず必要性を示す。$B_i$ はいずれも次数付き左右Noether環である。
実際 $R/(g)$ が次数付きNoetherで $g$ が中心なら、$(g)$ の随伴次数環は
$(R/(g))[t]$ の商であり、斉次イデアルの初項を有限生成して持ち上げられる。
内部次数を固定すれば正次数 $g$ の冪のフィルトレーションは有限なので、持上げは停止する。
有限な $B_4$ から逆向きに適用すればよい。
$B$ がAS正則なら、任意の有限生成次数付き $B_i$ 加群 $N$ について
$\underline{\operatorname{Ext}}^j_B(N,B)=0$ ($j<i$) が成り立つ。
$i=4$ では有限長とAS条件による。$i+1$ から $i$ へは
$N(-d)\xrightarrow{g_{i+1}}N$ の核・余核が有限生成 $B_{i+1}$ 加群であることを使う。
像で二つの短完全列に分ければ $j<i$ でこの写像はExtに同型を誘導する。
Extは有限生成自由解消の各項から内部次数が下に有界で、非零の次数移動と同型なら零となる。
以上は左でも右でも通る。$g_1,\ldots,g_{i-1}$ が既に正則列なら中心非零因子の二項分解を繰り返して
$\operatorname{Hom}_{B_{i-1}}(B_i,B_{i-1})$ は捩りを除いて $\operatorname{Ext}_B^{i-1}(B_i,B)$ に等しい。
これが零なので $g_i$ も非零因子である。四段の短完全列で第一条件が従い、第二条件は正則性による。

十分性では $B_4$ の頂点・次数別の語基底の正次数部分を $v$ と書き、各 $v$ を同じ経路で持ち上げる。
各 $B_i$ を順に割ってから持ち上げると、$B$ は順序つき記号
\[
                  g_1^{a_1}g_2^{a_2}g_3^{a_3}g_4^{a_4}v
\]
で張られる（末尾には頂点単位も許す）。実際、各段の核はその段での中心性により
$g_i B_{i-1}$ であり、正次数ゆえ次数帰納が使える。
端点に応じて対角元の対応成分だけを読む。この順序記号の個数は第一条件から各次数で標準値である。
第二条件により $3b$ 次までは既に独立である。

これを全次数へ延ばす。記号 $g_i,v$ の積に、$vv'$、$vg_i$、および $g_j g_i$ ($j>i$) を
それぞれ順序表示へ直す規則を入れる。$vg_i$ は $g_i v$ と、
より早い $g_s$ ($s<i$) が最初に立つ項の和とする。$g_jg_i$ も同様である。
この形は順次中心性と先ほどの持上げで常に選べる。$vv'$ は任意の順序表示とする。
内部次数固定のもとで、各書換えは陽に現れる $g_1,\ldots,g_4$ の個数の列を辞書順に増やすか、
同数なら末尾基底文字の個数を減らすか、どちらも同じなら逆転を一つ減らす。
$vv'$ から中心記号が出ない場合は一つの末尾基底文字だけになる。よって書換えは停止する。
規則は全て二記号に対するもので、重なりは三記号、重みは $3b$ 以下である。
この範囲の独立性により両側の簡約は同じ係数になり、ダイヤモンド補題で全順序記号が
これらの規則で定義した環の基底となる。
これが $B$ と同じ環であることも必要である。
元の矢を次数1の順序表示に送ると、二次関係は $2\le3b$ の独立性で零に簡約する。
また各 $g_i,v$ の元の経路式も重み $b$ 以下の独立性でその記号自身に戻る。
従って元の表示と規則による表示は互いに逆である。

最初の $g_1$ による左乗法は指数を一つ増やす。
その商でも次の $g_2$ について同じであり、順に中心なので両側の正則列を得る。
最後に各標準片側複体を $B_i$ に係数変更したものを $K_i$ とする。
正則列の短完全列により $K_i(-d_{i+1})\to K_i\to K_{i+1}$ は短完全列を与える。
$N=L+4=\sum d_i$ までの仮定から、$K_i$ の非増大複体のホモロジーはこの範囲で
内部次数が高々 $\sum_{j\le i}d_j$ でしか現れない。これは上の短完全列と $K_0$ の既知の消滅による帰納である。
$K_4$ 自体は $N$ より上に項を持たない。
$i=3,2,1,0$ の順に戻ると、$n>N$ では $H(K_i)_{n-d_{i+1}}\to H(K_i)_n$ は全射である。
最初に参照する $N$ 以下の次数は $\sum_{j\le i}d_j$ より大きいので零、以後は内部次数帰納で零となる。
よって $K_0$ は全次数で完全で、有限Frobenius対合が反対側とAS条件を与える。
\end{proof}
